\chapter{Fonctions linéaires et affines}
\label{corriges:09}
\section*{Exercice 9.1 -- Identifier et tracer}
$f$ est linéaire et affine ; $g$ est affine mais non linéaire ; $h$ est constante et affine, non linéaire ; $u$ n'est ni affine ni linéaire.
Pour $f$, placer $(0;0)$ et $(1;3)$. Pour $g$, placer $(0;4)$ et $(2;0)$. Relier chaque paire par une droite. Leur intersection vérifie $3x=-2x+4$, donc $(0{,}8;2{,}4)$.

\section*{Exercice 9.2 -- Retrouver une expression}
$a=(14-5)/(4-1)=3$. Comme $5=3\times1+b$, $b=2$, donc $f(x)=3x+2$. L'antécédent de $-1$ vérifie $3x+2=-1$, d'où $x=-1$.

\section*{Exercice 9.3 -- Interpréter les paramètres}
La température initiale est $22\degre$C. Le modèle prévoit une baisse de $0{,}4\degre$C par minute.
$T(15)=22-6=16\degre$C.
$22-0{,}4t=12$ donne $0{,}4t=10$, donc $t=25$ minutes, dans le domaine prévu. On ne prolonge pas automatiquement le modèle au-delà de $30$ minutes.

\section*{Exercice 9.4 -- Comparer deux abonnements}
$20+2n=5+3n$ donne $n=15$, coût $50$ euros. La différence vaut $A(n)-B(n)=15-n$. Elle est positive si $n<15$, donc $B$ est moins chère ; nulle pour $15$ ; négative si $n>15$, donc $A$ est moins chère à partir de $16$ utilisations.
